\section{引言}

本文关心一个最小但可审计的问题：当系统的“真实结构”是一个静态的约束系统 \(\Clo\)（而非一条先验给定的演化历史）时，观察者为何仍会产生丰富的“时空经验”？我们提出一个严格而可复现的答案：\textbf{时空来自观察者的向上搜索}；\textbf{时间来自为保持可逆性而必须携带的 uplift/历史胶带}。

\subsection{本体与观察者：两层区分}

在本体层（ontic），我们只固定一组一致性约束 \(\Clo\)。本体并不要求先计算出某条时间序列。相反，在观察层（epistemic），观察者只能访问投影后的可见量 \(y=\proj(x)\)，并在 \(\Clo\) 下尝试构造与之相容的提升（lifting）。此过程一般不是读取，而是受约束搜索；因此其代价天然具有实例依赖性，并可被解释为“观察者时间”的数学定义。

\begin{remark}[负信息时间与观察者时间]
本文把时间轴的本体成分具体化为 uplift 胶带：为使扩展态上的更新保持可逆，必须把“投影丢失的信息”以隐藏边标签/日志的形式记录下来。我们用其编码长度定义负信息时间 \(\tautime\)。同时，观察者时间 \(\kappaobs\) 仍以算法成本刻画“变慢/变复杂”。在最小假设下，二者在复杂度意义下同阶等价（\(\kappaobs=\Theta(\tautime)\)），从而把“时间=负信息”与“时间=计算代价”统一在同一可审计链条中。
\end{remark}

\subsection{Hilbert 尺子：把维度选择变成可计算输入}

观察者要做一致性搜索，必须选择一个“扫描尺子”：即如何组织变量的局部性、块结构与消元顺序。Hilbert 曲线族提供了一个同时具备\textbf{局部性保持}与\textbf{自相似递归}的统一工具：同一集合的自由度可在 1D/2D/3D 的 Hilbert 组织下产生不同的约束传播形态，从而改变搜索图结构与成本分布。进一步地，观察者可以在搜索过程中切换尺子（维度/尺度参数），把“遇到当前维度难以处理的问题就进入新维度扫描”的叙述翻译成一个严格的算法策略。

\subsection{与既有两条线索的关系}

本文框架不替代既有结果，而是把它们统一为“实例化的本体约束”：
\begin{itemize}
  \item 固定 \(m=6\) 的无损分块折叠核提供了“ontic 可逆 vs epistemic 搜索”的最小可计算模型；
  \item 一般 \(m\) 的 Zeckendorf 截断折叠与闭包图对象提供了“折叠边 + 微观几何边”在同一图上的统一组织方式。
\end{itemize}

